\chapter{Particle Lifecycle and Boundary Conditions}\label{ch:lifecycle}

A particle is not simply a position that updates each step: it carries a status that records whether it is
being advected, has been removed, has stranded, or has not yet been released. This chapter describes the four
statuses and the transitions between them (\cref{fig:lifecycle}), the record of why a particle died, the
treatments of the domain's boundaries --- removal at the lateral boundary, clamping at the surface and floor,
and the optional periodic fold --- the optional beaching model, and the delayed-release mechanism by which
particles enter the simulation on a schedule.

\section{The four statuses}\label{sec:life-status}

Each particle holds one of four statuses. An \code{ALIVE} particle is advected normally and is the only kind
included in the working set that the integrator processes. An \code{OUTSIDE} particle has left the domain; the
status is terminal, and the particle is thereafter inert, its position frozen where it last was valid. A
\code{BEACHED} particle has stranded in shallow water; uniquely among the non-alive statuses it is reversible,
as a beached particle may refloat (\cref{sec:life-beach}). A \code{NOT\_RELEASED} particle has not yet reached
its scheduled release time; it is excluded from the working set, frozen in place at no computational cost,
until it is activated (\cref{sec:life-release}). Frozen positions are copied from step to step unchanged, not
re-wrapped, so a dead particle at longitude $180^\circ$ stays there.

\begin{figure}[t]
\centering
\begin{tikzpicture}[node distance=14mm and 24mm]
  \node[statebox,fill=gray!12] (NR) {\code{NOT\_RELEASED}};
  \node[statebox,fill=laggreen!14,right=of NR] (AL) {\code{ALIVE}};
  \node[statebox,fill=kelp!12,right=of AL] (OUT) {\code{OUTSIDE}};
  \node[statebox,fill=lagelem!14,below=of AL] (BE) {\code{BEACHED}};
  \draw[trans] (NR)--(AL) node[midway,above,font=\footnotesize]{release};
  \draw[trans] (AL)--(OUT) node[midway,above,font=\footnotesize,align=center]{leaves domain\\(step reverted)};
  \draw[trans,transform canvas={xshift=-3pt}] (AL)--(BE) node[midway,left,font=\footnotesize]{beach};
  \draw[trans,transform canvas={xshift=3pt}] (BE)--(AL) node[midway,right,font=\footnotesize]{resuspend};
  \draw[trans,densely dashed] (OUT) |- node[pos=0.25,right,font=\footnotesize,align=left]{shallow or\\coastal anchor} (BE);
  \draw[trans] (AL) edge[loop above,looseness=5] node[font=\footnotesize,above,align=center]{surface/floor clamp;\\periodic fold} (AL);
\end{tikzpicture}
\caption{The particle status state machine. A particle is released from \code{NOT\_RELEASED} into
\code{ALIVE} at its scheduled time; it is removed to \code{OUTSIDE} if it leaves the domain (its step undone).
With beaching enabled, \code{ALIVE} and \code{BEACHED} interconvert, and a particle removed in the same step
with a shallow or coastal anchor is reclassified as \code{BEACHED} (dashed). Reaching the surface or floor
does not change the status: the depth is clamped and the particle remains \code{ALIVE}; nor does a fold back
into a periodic window (the self-loop).}
\label{fig:lifecycle}
\end{figure}

\section{Why a particle died}\label{sec:life-reason}

With \conf{mode.track\_failure\_reason: true} each particle also carries a failure reason: none,
\emph{interpolation failed} (a stage found no fit or no containing element) or \emph{out of bounds} (the final
position lies in no element). Two further codes, for surface and bottom contact, are reserved and never set,
because those limits never remove a particle (\cref{sec:life-bc}). The reason is set once, when the particle
dies, and never overwritten; every integration path records it, fused and unfused alike (\cref{ch:fused}). A
beaching remap or a resuspension clears it, so the field keeps the meaning ``cause of death''. It is written
to the trajectory output and to checkpoints.

\section{Lateral removal versus vertical clamping}\label{sec:life-bc}

The engine treats the lateral boundary and the vertical limits in qualitatively different ways, and the
distinction is deliberate. The containing-element test of \cref{sec:loc-contain} is the lateral boundary
condition: a particle whose advected position is contained by no element has left the domain, and the engine
\emph{removes} it --- it sets the status to \code{OUTSIDE} and reverts the position to the start of the step,
so the particle does not record a spurious move across the boundary. A failed interpolation removes a particle
in the same way, and so does a diffusive kick that ends outside the mesh (\cref{sec:diff-composition}).
Because the test is geometric, an open boundary and a coastline are indistinguishable to it; both produce an
\code{OUTSIDE} result unless the beaching model says otherwise. The surface and the sea floor, by contrast,
are enforced by \emph{clamping} after the advective step (\cref{sec:vert-clamp}): a particle that would cross
either limit is held at it and continues to be advected. The vertical diffusive kick is the exception: it
\emph{reflects} at the limits, because a clamp would pile particles up there (\cref{sec:diff-vertical}).

\section{Periodic windows}\label{sec:life-periodic}

An idealized channel domain is made periodic in longitude, latitude or both by giving a window
$[x_0,x_1)$ in \conf{numerics.periodic\_lon} or \conf{numerics.periodic\_lat}. After every step, before the beaching hook and the
record, an alive particle outside the window is folded back,
\begin{equation}\label{eq:life-fold}
  x' = x_0 + \bigl((x-x_0)\bmod(x_1-x_0)\bigr),
\end{equation}
and receives a fresh nearest node; particles inside the window are left bit for bit untouched. Seeds are
folded once at the start, including those not yet released. The fold adds no connectivity across the seam, so
the mesh must extend beyond each periodic side by a \emph{ghost band} at least one step's displacement wide:
a particle must still be inside the mesh at the end of the step in which it crosses. The configuration is
checked before the output file is created: the windows must lie strictly inside $(-180^\circ,180^\circ)$ and
$(-90^\circ,90^\circ)$ and overlap the mesh, a periodic-longitude mesh must not straddle the $\pm180^\circ$
meridian, and a missing ghost band on either side produces a warning. Periodic windows cannot be combined with
a rotated frame (\cref{ch:frame}). The windows are written to the output metadata when set.

\section{Beaching and resuspension}\label{sec:life-beach}

When beaching is disabled (the default), every lateral exit becomes \code{OUTSIDE}. When the key
\conf{numerics.beaching.enabled} is set, a post-step hook strands particles by up to four criteria, combined
inclusively. The water column at the particle's nearest node is taken raw, $H=d+\eta$ when
\conf{use\_eta} is set and $H=d$ otherwise, without the minimum-thickness floor, which would hide tidal
drying; $\eta$ comes from the step's own velocity payload.
\begin{description}[font=\normalfont\conf,style=nextline,leftmargin=1.5em]
  \item[depth] $H<$ \conf{depth\_min}, or the node is marked dry by \conf{dry\_mask\_var}.
  \item[coastal\_element] the particle's containing element has an edge on the coast (an edge owned by one
    element only).
  \item[on\_failure] a particle the integrator has just removed is reclassified as \code{BEACHED} if its
    reverted anchor has $H$ below the zone depth or lies in a coastal element (the coastal test needs the
    coastal flags, built only when \conf{coastal\_element} or \conf{stochastic} is also a criterion); a removal in deep water at an
    open boundary stays \code{OUTSIDE}.
  \item[stochastic] within a near-shore zone --- $H<$ \conf{zone\_depth} (default twice \conf{depth\_min}), or
    coastal contact whenever the coastal flags exist --- a particle strands with probability
    $1-\exp(-\lvert\dt\rvert/t_{\text{beach}})$ per step.
\end{description}
A beached particle counts as re-wetted when $H>$ \conf{depth\_min} $+$ \conf{rewet\_hysteresis} and its node is
not marked dry; the hysteresis suppresses flicker between the two states. The resuspension policy is one of
\conf{never}; \conf{auto}, which refloats a particle after \conf{min\_wet\_steps} consecutive re-wetted steps;
\conf{stochastic}, which refloats a re-wetted particle with probability $1-\exp(-\lvert\dt\rvert/t_{\text{resuspend}})$;
and \conf{auto\_or\_stochastic}. The random draws use the counter-based generator of \cref{sec:diff-rng} with
their own components, keyed by the original particle index and the global step, and the count of re-wetted
steps is written to checkpoints, so a resumed run reproduces the uninterrupted one bit for bit.

The hook runs in two phases after the step: first the reclassification of just-removed particles, then the
criteria and the resuspension. Its containment test runs only for alive, beached and just-removed particles. The
coastal flags are built only when a criterion needs them. A configuration asking for $\eta$ without a
sea-surface height in the dataset, or naming a dry-mask variable that is absent, is refused at startup. On a
HEALPix mesh with neither a bathymetry variable nor \conf{mesh.levels} the node depths are zero, and the depth criterion would strand every particle.

\section{Delayed and pulsed release}\label{sec:life-release}

Release positions are read from the seed file. Two library helpers can write one: a regular grid over the
mesh extent, shrunk inward by a margin and filtered to points inside the mesh, and points scattered uniformly
in a circle of given radius around a centre, scaled by $\cos\latx$ so that the circle is round in metres.
Both work in the mesh's own coordinates: on a rotated mesh the caller rotates the centre into the mesh
frame and the kept points back to geographic before saving them as seeds.

Particles need not all enter at the start of the run. A per-particle release offset (a fourth column of the
seed file, in seconds from the start) and a pulse schedule that re-releases the whole seed set at a fixed
interval (\conf{time.n\_pulses}, \conf{time.pulse\_interval}) are both supported, and both are implemented by
pre-allocation and freezing rather than by growing the particle array. At setup the seed set is replicated to
its full final size --- the number of pulses times the number of seeds, pulse by pulse --- and each copy is given
a release step
\begin{equation}\label{eq:life-release}
  k_{\text{rel}}=\operatorname{round}\Bigl(\frac{t_{\text{rel}}+p\,\Delta t_{\text{pulse}}}{\dt}\Bigr)
\end{equation}
for pulse $p$, with a warning when the offset is not a whole number of steps. Particles whose release step
lies in the future begin as \code{NOT\_RELEASED}. At the top of each step the engine activates those whose
step has come, $k\ge k_{\text{rel}}$ in the global step count (so the rule holds for backward runs and across
restarts), flipping them to \code{ALIVE} so they are advected during that same step. A frozen particle costs
nothing in the kernel, which skips it, and the fixed array size is what allows the streaming output writer to
use a constant trajectory dimension. Dormant rows are written as missing values with their status; a particle
whose release step lies beyond the end of the run produces a warning and stays missing throughout. The output
records each particle's release position and release time, the start of the step in which it is activated
(\cref{eq:out-release}), so the rounded schedule can be read back from the file. The
diffusive noise is keyed by the global step, so a particle released later draws a different sequence from one
released at the start, by design. A run with a single release and no offsets reduces exactly to the
unscheduled case.

\codeanchor{The statuses and failure reasons are the enumerations \fn{ParticleStatus} and \fn{FailureReason} in
\srcfile{core/particles.py}; lateral removal and the revert are in the integrator kernels of
\srcfile{numerics/integrators.py}; the periodic fold and its checks are \fn{\_apply\_periodic\_wrap} and
\fn{\_validate\_periodic\_windows}, the beaching hook is \fn{\_apply\_beaching\_hook} (\srcfile{tracking.py})
over \fn{evaluate\_beaching} (\srcfile{numerics/beaching.py}); the release schedule is
\fn{build\_release\_state}, \fn{activate\_released} and \fn{mask\_dormant\_for\_record}
(\srcfile{core/release.py}).}
